\section{Eserciziario Tematico Selezionato dalle Prove Scritte}
\label{sec:eserciziario_tematico}

In questa sezione vengono raccolti esercizi emblematici tratti dagli esempi d'esame storici, suddivisi per nuclei concettuali per facilitare il ripasso mirato.

\begin{esercizio}{Probabilità Bivariata e Regressione Lineare}{es_tematico_bivariata}
Sia $(X, Y)$ un vettore aleatorio continuo con densità congiunta $f_{X,Y}(x,y) = \frac{6}{5}(x + y^2)$ sul quadrato unitario $[0,1]^2$.
\begin{enumerate}
    \item[(a)] Determinare le densità marginali $f_X(x)$ e $f_Y(y)$.
    \item[(b)] Calcolare la covarianza $\operatorname{Cov}(X, Y)$ e verificare se $X$ e $Y$ sono incorrelate.
\end{enumerate}
\end{esercizio}

\begin{dimostrazione}
\textbf{(a) Densità Marginali su $[0,1]$:}
\begin{align*}
    f_X(x) &= \int_0^1 \frac{6}{5}(x + y^2) \, dy = \frac{6}{5} \left[ x y + \frac{y^3}{3} \right]_0^1 = \bm{\frac{6}{5}\left(x + \frac{1}{3}\right), \qquad x \in [0,1]} \\
    f_Y(y) &= \int_0^1 \frac{6}{5}(x + y^2) \, dx = \frac{6}{5} \left[ \frac{x^2}{2} + x y^2 \right]_0^1 = \bm{\frac{6}{5}\left(\frac{1}{2} + y^2\right), \qquad y \in [0,1]}
\end{align*}

\textbf{(b) Momenti e Covarianza:}
\begin{align*}
    \mathbb{E}[X] &= \frac{6}{5} \int_0^1 x\left(x + \frac{1}{3}\right) dx = \frac{6}{5}\left[\frac{1}{3} + \frac{1}{6}\right] = \frac{6}{5} \cdot \frac{1}{2} = \bm{\frac{3}{5}} = 0.6 \\
    \mathbb{E}[Y] &= \frac{6}{5} \int_0^1 y\left(\frac{1}{2} + y^2\right) dy = \frac{6}{5}\left[\frac{1}{4} + \frac{1}{4}\right] = \frac{6}{5} \cdot \frac{1}{2} = \bm{\frac{3}{5}} = 0.6 \\
    \mathbb{E}[XY] &= \frac{6}{5} \int_0^1 \int_0^1 (x y)(x + y^2) \, dx \, dy = \frac{6}{5} \int_0^1 y \left[ \frac{x^3}{3} + \frac{x^2 y^2}{2} \right]_0^1 dy = \frac{6}{5} \int_0^1 \left(\frac{y}{3} + \frac{y^3}{2}\right) dy \\
    &= \frac{6}{5} \left[ \frac{1}{6} + \frac{1}{8} \right] = \frac{6}{5} \cdot \frac{7}{24} = \bm{\frac{7}{20}} = 0.35
\end{align*}
La covarianza è:
\begin{equation*}
    \operatorname{Cov}(X, Y) = \mathbb{E}[XY] - \mathbb{E}[X]\mathbb{E}[Y] = 0.35 - (0.6)(0.6) = 0.35 - 0.36 = \bm{-0.01} \neq 0
\end{equation*}
Le variabili presentano una lieve covarianza negativa e dunque \textbf{non sono incorrelate} (e di conseguenza non sono indipendenti). $\blacksquare$
\end{dimostrazione}
